\chapter{电池把化学变化变成电流}

\begin{kaichang}
电视遥控器又不灵了。你抠出那两节五号电池，在手心里搓了搓、捂了捂，再装回去——嘿，又能换台了。不过撑不了几天，它又会彻底罢工。

先别急着扔。想一想这几个问题：电池上印着“\ud{1.5}{V}”，这个“电”是从哪儿来的？是出厂的时候灌进去的吗？“电池用完了”，用完的到底是什么——是里面的“电”漏光了，还是别的东西耗尽了？还有，为什么搓一搓、捂一捂，它又能回光返照一小会儿？

这一章结束时，你会看到：一节普通电池里发生的事情，和你呼吸时每个细胞里发生的事情，原理竟然是同一套。
\end{kaichang}

\section{把电池拆开之前，先看看它会什么}

不拆电池（碱性电池内部有腐蚀性物质，交给回收站，别自己动手拆），只看它对外表现出的行为，就能列出一串现象：

\begin{itemize}[itemsep=2pt]
\item 它有两个明确的端子：一个标“$+$”，一个标“$-$”。接反了，有些电器就不工作。
\item 接上小灯泡或电动机，电流就持续流动，灯泡能亮很久——不是静电那种“啪”一下的火花，而是稳定持久的输出。
\item 用久了，灯泡变暗、电动机变慢，最后停转。电池“死了”。
\item 死去的电池并没有变轻多少，外壳也还是那个外壳。它失去的不是某种看得见的东西。
\item 一节新干电池正负极之间的电压约为 \ud{1.5}{V}；把两节首尾相接地串起来，电压变成约 \ud{3}{V}；并联起来，电压不变，但能供电更久。
\end{itemize}

还有一件更好玩的事：根本不用买电池。把一个柠檬切半，往果肉里插一枚镀锌铁钉和一段铜线，用导线把铁钉和铜线连到一颗发光二极管（LED）上——二极管居然微微发亮。柠檬又不是工厂造的，它的“电”从哪来？

这些现象拼在一起指向一个结论：电池输出的不是神秘流体，而是\textbf{一场正在进行的化学反应}。反应物耗尽了，电池就死了。要验证这个想法，得把镜头推进到粒子层。

\section{把氧化还原“拆开用”}

第 33 章讲过电子转移推动的反应：锌片插进硫酸铜溶液里，锌原子把电子直接交给铜离子，锌变成锌离子溶解掉，铜离子变成铜原子析出来，溶液的蓝色慢慢褪去，锌片表面裹上一层红铜，同时放热。电子从锌直接跳到铜，能量全部变成热，一点有用的功都没留下。

现在想一个朴素的问题：能不能逼着电子\textbf{绕一段远路}？让锌在一个房间交出电子，让铜离子在另一个房间接收电子，中间用一根导线连接——电子要过去，只能穿过导线，顺便推动灯泡里的灯丝发光。这个“让氧化和还原分家”的装置，就是\textbf{原电池}。

\begin{figure}[htbp]
\centering
\begin{tikzpicture}[scale=0.9]
% 左烧杯
\draw (0,0) -- (0,2.2) (2,0) -- (2,2.2) (0,0) -- (2,0);
\node at (1,-0.4) {锌离子溶液};
\draw[thick] (0.8,1.1) -- (1.2,1.1) (0.8,2.6) -- (1.2,2.6) (0.8,1.1) -- (0.8,2.6) (1.2,1.1) -- (1.2,2.6);
\node at (1,3.0) {锌片($-$)};
% 右烧杯
\draw (5,0) -- (5,2.2) (7,0) -- (7,2.2) (5,0) -- (7,0);
\node at (6,-0.4) {铜离子溶液};
\draw[thick] (5.8,1.1) -- (6.2,1.1) (5.8,2.6) -- (6.2,2.6) (5.8,1.1) -- (5.8,2.6) (6.2,1.1) -- (6.2,2.6);
\node at (6,3.0) {铜片($+$)};
% 外电路
\draw (1,2.6) -- (1,3.8) -- (6,3.8) -- (6,2.6);
\draw[->,thick] (2.5,3.8) -- (4.5,3.8);
\node at (3.5,4.2) {电子经导线流向铜片};
% 盐桥
\draw[thick] (1.6,1.6) .. controls (2.5,3.0) and (4.5,3.0) .. (5.4,1.6);
\node at (3.5,2.6) {盐桥(离子通道)};
\end{tikzpicture}
\caption{铜锌原电池的示意画法。这是人为的表示法：烧杯大小、盐桥形状都不代表真实比例。}
\end{figure}

装置工作时，三个动作同时进行，缺一不可：

\begin{itemize}[itemsep=2pt]
\item \textbf{锌的一极}：锌原子不断交出电子，变成锌离子进入溶液。这一极叫\textbf{负极}（电子从这里出发），发生的反应叫氧化。
\item \textbf{铜的一极}：电子顺着导线到达铜片表面，把溶液里凑过来的铜离子“签收”成铜原子，镀在铜片上。这一极叫\textbf{正极}，发生的反应叫还原。
\item \textbf{内部}：锌那边正电荷越积越多，铜那边正电荷越来越少，如果不平衡，电子流几步就会“堵车”。盐桥（一条灌满惰性盐溶液的管子）让离子在里面迁移：阴离子移向锌侧，阳离子移向铜侧，把两边的电荷账目抹平。
\end{itemize}

请注意一个容易想歪的地方：\textbf{电子不走溶液，离子不走导线}。外电路里是电子在金属中移动，电池内部是离子在液体中迁移，两条“传送带”在电极表面交接，构成一个完整的环。任何一段断开——拔掉导线，或者拿走盐桥——整个环立刻停摆，反应随即冻结。这正是为什么一节外皮破损、电解液干涸的旧电池会死得透透的。

再追一层：电子为什么不能干脆穿过溶液抄近道？因为自由电子在金属里才跑得动，溶进水里立刻被水分子“抓获”，寸步难行；溶液导电靠的是离子整体搬家，不是电子飞奔。所以只要装置搭得对，电子别无选择，只能老老实实走外电路。反过来，如果你用一根导线把电池正负极\textbf{直接}短接，电子倒是畅通无阻了，可落差全部兑换成了导线和电池内部的热量——这就是“短路会发烫甚至起火”的原因，也是第 27 章能量账的又一次提醒：自由能落差总要花掉，区别只在花成有用功还是花成热。

\section{电压是“电子的海拔差”}

为什么电子愿意从锌流向铜，而不是反过来？第 28 章的自由能图景在这里照样适用：锌失电子、铜离子得电子这一套组合，总自由能是下降的，所以它能自发进行。两种金属“交出电子的难易程度”不同，就像两个高度不同的水面——接通之后，水往低处流，电子往“能量低处”流。这个落差，就是电压表上读到的\textbf{电势差}。

换一对金属，落差就变了。锌和铜组合，约 \ud{1.1}{V}；换成镁和铜，落差更大，约 \ud{2.7}{V}。把两节电池串联，相当于把两段瀑布叠起来，落差相加——这就是串联升压的原因。所以电池上印的电压不是“电量”，而是这场化学反应的\textbf{落差大小}；电池能撑多久，才取决于里面存了多少反应物。

\begin{qiaoliang}
一节碱性五号电池大约含 \ud{3}{g} 金属锌，能放出多少电荷？算一笔账：锌的摩尔质量约为 \ud{65}{g\cdot mol^{-1}}，\ud{3}{g} 锌约是 $3/65\approx 0.046$ mol；每个锌原子交出 2 个电子，共约 0.092 mol 电子。1 mol 电子携带的电荷约为 \ud{96500}{C}（这个常数叫法拉第常数），所以总电荷约 $0.092\times 96500\approx 8900$ C。换算成电池行业爱用的单位：$8900/3600$ 约等于 \ud{2.5}{A\cdot h}——和市售五号碱性电池标称的容量正好同一个数量级。换句话说，电池容量根本不需要背：数数里面有多少反应物、每个反应物交几个电子，账就出来了。
\end{qiaoliang}

到这一步，可以请符号登场了。两个“半反应”分开写：

\begin{itemize}[itemsep=2pt]
\item 负极（氧化）：\[\chem{Zn} \rightarrow \chem{Zn^{2+}} + \chem{2e^-}\]
\item 正极（还原）：\[\chem{Cu^{2+}} + \chem{2e^-} \rightarrow \chem{Cu}\]
\end{itemize}

加起来，电子在等式两边抵消，就是第 33 章见过的总反应 $\chem{Zn} + \chem{Cu^{2+}} \rightarrow \chem{Zn^{2+}} + \chem{Cu}$。同一本账，只是这一回电子被要求“走外路”，把自由能的落差兑换成了电流。书写电池时还有一套极简记法：$\chem{Zn}\,|\,\chem{Zn^{2+}}\,\|\,\chem{Cu^{2+}}\,|\,\chem{Cu}$——单竖线表示相界面，双竖线表示盐桥，左边负极，右边正极。一行符号，就把整个装置说完了。

\section{青蛙腿引发的一场科学争论}

\begin{zhengju}
电池的诞生，始于一次“冤案”。18 世纪 80 年代，意大利解剖学家伽伐尼发现：剥好的青蛙腿挂在铜钩上，一碰到铁栏杆就猛烈抽搐，像复活了一样。当时静电实验正时髦，伽伐尼给出的解释是：青蛙体内储藏者“动物电”，金属只是把它放了出来。他还做了对照——只用一种金属时抽搐弱得多，两种不同金属接触时最强烈，于是他更确信：电是青蛙的，金属只是引线。

他的同胞伏打读到了论文，却盯上了另一条线索：为什么一定要两种\textbf{不同}的金属？伏打怀疑电根本不是青蛙的，青蛙腿只是一台灵敏的“验电器”（蛙腿含盐水，既是导体又能抽动显示电流）。要证伪伽伐尼，就得彻底抛开青蛙。伏打把银片和锌片交替叠起来，每两片之间垫一块浸满盐水的布，叠了几十层——1800 年，这台“伏打电堆”输出了持续稳定的电流，不需要任何青蛙。这是人类历史上第一个能持续供电的电源，物理学从此有了研究稳定电流的工具，十几年后电解水、电解出钾钠等元素都靠它。

谁赢了？各对一半。伏打对了：电堆的电来自金属与盐水之间的化学反应，不来自动物。但伽伐尼也没全错：几十年后人们确认，神经和肌肉真的靠电信号工作（你的心脏此刻就受电信号指挥），只是那套“生物电”的机制要到很久以后才弄清楚——它同样建立在本章的离子迁移原理上，第 53 章还会遇到它的升级版。这场争论的方法值得记住：伏打没有停留在“我觉得不对”，而是设计了一个把可疑因素（青蛙）整个剔除的实验，让变量只剩金属和盐。
\end{zhengju}

\section{从柠檬到锂电池：电池家族巡礼}

铜锌电池是教学模型，真实世界的电池是同一个原理的四种工程变体：

\begin{itemize}[itemsep=2pt]
\item \textbf{锌锰干电池}（普通一次性电池）：锌皮外壳直接当负极，中央的碳棒是正极的“集电棒”，正极材料是二氧化锰，电解质是糊状的氯化铵或碱液。把液体做成“糊”，是为了让你能把电池横着放、倒着放——工程上这叫防漏，化学上本质没变。
\item \textbf{铅酸电池}（汽车电瓶）：负极是铅，正极是二氧化铅，电解质是稀硫酸，单格约 \ud{2}{V}，汽车电瓶把 6 格串联成 \ud{12}{V}。它又重又丑，却有一个绝技：能瞬间输出几百安培的大电流去带动启动马达，而且便宜、耐造、回收体系成熟，所以一百多年没被淘汰。
\item \textbf{锂离子电池}（手机、电动车）：负极用能“嵌”锂的石墨，正极用钴酸锂或磷酸铁锂等材料。它的聪明之处是把锂做成\textbf{离子}在两层材料之间来回搬家：放电时锂离子从负极搬家到正极，充电时搬回来，两边材料的骨架基本不动——所以它能反复充放上千次。锂是最轻的金属，又特别爱丢电子，单位重量存下的能量是铅酸电池的好几倍，这就是手机能越做越薄的化学原因。
\item \textbf{燃料电池}：前面的电池都把反应物封在壳子里，燃料电池却是“敞开吃”——氢气（或天然气重整出的氢）和氧气持续通进去，在第 30 章说的催化剂表面分别交出和接受电子，电子被迫走外电路，产物只有水。它更像一座微型发电站而不是储能罐：只要燃料不断，电就不断。
\end{itemize}

锂电池能量密度高，代价是“脾气”也更急。它的电解质是可燃的有机液体，一旦内部短路（摔伤、过充、生产缺陷都可能引发），短路点会急剧发热，热量又加速副反应，副反应再放热——这个自我加速的恶性循环叫“热失控”，最后整节电池喷火。这就是为什么民航限制充电宝托运、为什么鼓包的手机电池必须立刻停用。工程师的对策也全是化学与材料手段：更耐高温的隔膜、添加阻燃剂、用磷酸铁锂这类放热更温和的正极材料。安全不是一句“小心使用”，而是一连串在分子层面做的权衡。

一个共同的判断标准浮现出来：选哪种电池，看的不是谁“先进”，而是谁的\textbf{反应物便宜、可逆性好、单位重量的自由能落差大、副反应少}。电池工业一百多年的竞争，就是在这几条之间做取舍。

还有一笔账要算在电池“死后”。废电池里有铅、镉、钴、镍等金属离子，随手扔掉的电池在垃圾填埋场里慢慢腐烂，金属离子随雨水渗进土壤和地下水，最终可能回到食物链里——重金属离子爱在生物体内富集，干扰酶和神经系统。但换个角度看，这些“污染物”其实是\textbf{浓度很高的矿石}：一吨废旧锂电池里的钴和锂，比一吨天然矿石里的含量还高。所以电池回收不只是环保口号，还是一门正经的城市采矿工业：把废电池拆解、溶解，再用本章讲的电解和第 35 章的沉淀、配位方法把各种金属离子一样一样分离提纯出来。开头说旧电池要交给回收站，现在你知道那不只是为了不弄脏土地，也是为了把“落差”的原材料接回产业链。

\section{充电不是“把电子塞回去”}

可充电电池放电时，反应自发地朝一个方向跑；充电时，外接充电器强行施加一个反方向的更大电压，\textbf{硬把电子推回去}，让反应逆转——正极变回二氧化铅、锂离子搬回石墨。所以充电消耗的是充电器送来的电，电池里既没有“电被灌进去”，也没有“电被用完”，有的只是化学物质在两个方向之间来回倒腾。能量账在第 27 章就立过规矩：充进去的电，一部分变成化学物质的自由能存起来，一部分不可避免地变成热散掉，永远存不满、也取不尽。

那为什么不是所有电池都能充电？一次性碱性电池里，放电产物会形成难以逆转的新物质（比如锌变成氧化锌粉末，形状、位置都变了），强行反充不但恢复不了原来的结构，还可能析出气体、鼓包甚至爆裂——这就是为什么充电器说明书上写着“严禁给一次性电池充电”，这条禁令背后是实打实的化学。

即便是锂电池，每次搬家也做不到百分百完美：极少数锂离子会“迷路”，和电解质发生副反应、在负极表面析出金属锂，可用容量便一点点缩水。你的手机电池用两年后续航明显变短，不是电“漏”了，而是能继续搬家的锂离子变少了。低温下离子搬得慢、大电流下副反应多，所以冬天电动车续航缩水、快充伤电池——全是动力学问题（第 29 章）。

\begin{shiyan}
\textbf{水果电池}（安全，可在家做）。

材料：一个柠檬（或土豆、苹果）、一枚镀锌铁钉（表面有锌）、一段粗铜线或一枚黄铜钥匙、几根带鳄鱼夹的导线、一颗低电压 LED（可选：万用表）。

步骤：①把铁钉和铜线并排插进果肉，相距约 \ud{2}{cm}，注意二者在果肉里\textbf{不要碰在一起}；②用导线把铁钉和铜线分别连到万用表上，读电压；③换不同的水果、不同的金属组合（比如铁钉配铁钉），记录电压变化；④若有 LED，把 3～4 个水果电池串联起来再试，观察能否点亮。

观察点：锌和铜组合约能读到 \ud{0.8}{V}～\ud{1.0}{V}；换成两枚相同的铁钉，电压几乎归零——亲自验证“落差来自两种不同的金属”。

安全提示：实验本身电压极低，没有触电危险；但\textbf{绝对不要}把任何自制电池接到墙壁插座上，实验后的水果已被金属离子污染，请丢弃不要食用，做完洗手。
\end{shiyan}

\section{把电反过来用：电解、电镀与炼铝}

原电池是“自发反应推着电子走”，把电源反接、强行驱动一个本来不自发的反应，就叫\textbf{电解}。同一个原理翻个面，撑起了一整条工业链：

\begin{itemize}[itemsep=2pt]
\item \textbf{电镀}：把铁钉挂到含铜离子（或镍、铬离子）的溶液里当负极，通直流电，铜离子就被强行按在铁钉表面还原成致密的金属层——防锈又美观。你的水龙头、钥匙、眼镜框多半都镀过。
\item \textbf{电解精炼铝}：铝太爱丢电子，自然界只以氧化物存在，木炭还原法拿它没办法。工业上把氧化铝溶进熔融的冰晶石里（纯氧化铝熔点超过 $2000\,^{\circ}\mathrm{C}$，冰晶石能把它降到约 $960\,^{\circ}\mathrm{C}$，省下巨额电费），再用强大的电流把铝离子硬还原成液态铝。铝曾经是比银子还贵的“贵族金属”，电解法一成熟，它才变成做易拉罐的大路货——你手里那罐饮料，是电化学工业平权运动的成果。
\item \textbf{氯碱工业}：电解饱和食盐水，同时得到三样宝贝——烧碱（\chem{NaOH}）、氯气（\chem{Cl_2}）和氢气（\chem{H_2}）。肥皂、纸张、消毒剂、塑料原料 \chem{PVC} 的背后都站着这条生产线。
\end{itemize}

注意方向感：电解池里，接电源正极的那极发生氧化（叫阳极），接负极的那极发生还原（叫阴极）——和原电池里“负氧正还”的说法看似拧着，其实“氧化总在阳极、还原总在阴极”是统一的，只是电极的正负号随装置角色翻转。记不住的诀窍是：盯住电子往哪边走，一切称呼都能推出来，不用背。

\section{这个模型什么时候会失灵}

本章的图景——两个半反应、电子走外电路、电压是落差——在解释“电池为什么能供电”时非常可靠，但有几处它是简化过的：

\begin{itemize}[itemsep=2pt]
\item \textbf{电压不是常数}。标称 \ud{1.5}{V} 只是新电池、几乎不带负载时的读数。反应物消耗、产物堆积，落差会一路下滑；接上耗电大的电器，电池内部离子迁移跟不上（内阻），电压还会当场下跌。旧电池“搓一捂又能用”，正是利用了这一点：温度升高让离子迁移变快（第 29 章），电压短暂回升，但反应物总量一点没变多——回光返照，不是复活。
\item \textbf{“半反应”是记账拆分，不是先后两步}。锌放电子和铜得电子在时间上同步发生、速率精确匹配，谁先谁后的问题没有意义。
\item \textbf{电压读不出电量}。一颗快没电的锂电池静置时电压仍接近满电，一上大电流就瞬间崩盘。判断余量要靠放电曲线的积分（电量计芯片干的活），不是瞄一眼电压。
\item 浓差本身也能造出电压（同一物质两边浓度不同就有电势差，叫浓差电池），说明“落差”比“两种金属”更本质——这正是挑战层能斯特方程要说的。
\end{itemize}

\begin{wuqu}
“电池里存着很多电子，用电器把电子用光了，电池就没电了。”——错得很有迷惑性，因为日常语言就这么说。反例随手可举：第一，电流是环流，从电池流出的电子数和流回的电子数\textbf{严格相等}，一个也不会少，用电器消耗的不是电子本身，而是电子从高电势落到低电势时释放的能量，正如水车消耗的不是水而是落差；第二，一节死去的电池里电子一个不少（锌离子、氧化产物里都满满是电子），耗尽的是\textbf{能够自发进行反应的反应物}——也就是那份自由能落差。把“电池存电”改口成“电池存着一份还没花掉的化学落差”，你对所有电池新闻（容量、快充、衰减）的判断力都会立刻升级。
\end{wuqu}

\section{回到遥控器里那两节电池}

现在兑现开场的承诺。“\ud{1.5}{V}”不是灌进去的电，而是锌与二氧化锰之间那场反应的自由能落差；“电池用完”，用完的不是电子，而是能交电子的锌和能收电子的二氧化锰——落差抹平了，反应抵达平衡，电流也就断了；至于搓一搓、捂一捂的“复活术”，只是加热加速了离子迁移、让残余反应物接触得充分一点，落差短暂显形。它改变不了总账，所以注定撑不久——能量守恒这本大账，谁也赖不掉。

顺手还能解开柠檬之谜：柠檬汁里的酸提供氢离子（第 32 章），锌钉比铜线更爱丢电子，果肉就是天然盐桥。你的水果电池和伏打 1800 年的电堆，是同一件东西。

\section{你身体里有亿万块“电池”}

这一章讲的是把化学反应的自由能变成电流。如果把镜头转向生物体，你会发现生命把这个把戏玩到了极致：你的每个细胞都在膜两侧维持着离子浓度差——里面钾多、外面钠多，就像一块永远处于“半充电”状态的微型电池；神经冲动是这道电势差沿细胞膜飞快翻转的波浪；而你吃进去的葡萄糖，最终正是在线粒体里通过一条“电子传送链”，让电子一级级下落、用落差泵出质子梯度来合成 ATP。电池的原理没有变，只是把锌和铜换成了蛋白质和离子。第 51 章讲细胞膜、第 53 章讲细胞呼吸时，我们会带着今天的“半反应、电势差、离子迁移”三件套回来——到那时你会看到，线粒体其实是一座比你手机里那块锂电池精巧得多的燃料电池。

\begin{tiaozhan}
（跳过不影响主线）怎样定量计算“落差”？化学家给每种半反应测了一个\textbf{标准电极电势} $E^{\ominus}$（相对于氢电极，单位 V），查表相减即得电池电压：$E^{\ominus}_{\text{池}} = E^{\ominus}_{\text{正}} - E^{\ominus}_{\text{负}}$。铜锌电池：$0.34-(-0.76)$，即 \ud{1.10}{V}。$E^{\ominus}$ 与自由能只差一个换算：$\Delta G^{\ominus}=-nFE^{\ominus}$，$n$ 是转移电子数，$F$ 是法拉第常数。浓度不是标准状态时，用\textbf{能斯特方程}修正：$E = E^{\ominus} - \dfrac{RT}{nF}\ln Q$，$Q$ 是产物与反应物的活度比。这个方程同时解释了三件事：电池越用电压越低（产物 $Q$ 变大）；浓差电池的存在（两边 $Q$ 不同就有 $E$）；以及第 31 章平衡常数与电压的换算（平衡时 $E=0$，$Q=K$）。电化学和热力学，在这个方程里握了手。
\end{tiaozhan}

\section*{练一练}

\begin{sikaoti}
\item 画一幅铜锌原电池的粒子图：两个烧杯里分别画锌离子、铜离子和硫酸根离子，标出电子在导线中的方向、离子在盐桥中的迁移方向。图旁注明：这是人为表示法，离子并不长这样。
\item 水果电池实验中，把“锌钉$+$铜线”换成“两枚相同的锌钉”，电压几乎为零。用“自由能落差”解释为什么；再预测：换成“镁条$+$铜线”，电压会怎样变化？
\item 一节电池在电路中工作了一小时。流入用电器的电子数与流出用电器的电子数哪个多？用电器究竟“用掉”了什么？（提示：想想水车和水。）
\item 给铅酸电池充电时，外接电源扮演了原电池中谁的角色？画一张能量流图，标出“充电器电能→化学物质自由能→散掉的热”三条支流。
\item 电解食盐水与铜锌原电池用的是同一套“电子转移”语言。请各写一句话说明：在两个装置里，分别是什么力量在推动电子移动？这两种力量有何本质不同？
\end{sikaoti}
